\documentclass[a4paper]{article}
\usepackage{aqis}
\usepackage{enumitem}
\usepackage{graphicx}
\usepackage{array}
\usepackage{booktabs}
\usepackage{multirow}
\usepackage{tikz}
\usetikzlibrary{calc,arrows.meta}
\usepackage[colorlinks]{hyperref}
\usepackage{color}
\newcommand{\blue}[1]{\textcolor{blue}{#1}}

\newcommand{\rd}{r}
\newcommand{\gr}{g}
\newcommand{\bl}{b}
\newcommand{\GF}{\mathrm{GF}(2)}
\providecommand{\ket}[1]{\ensuremath{|#1\rangle}}
\providecommand{\bra}[1]{\ensuremath{\langle#1|}}
\newcommand{\eh}[1]{e^{\rightarrow}_{#1}}
\newcommand{\ev}[1]{e^{\uparrow}_{#1}}
\newcommand{\ed}[1]{e^{\nearrow}_{#1}}
\newcommand{\cs}[1]{\texttt{#1}}
\newcommand{\Pch}{P_{\mathrm{charge}}}
\newcommand{\Pfl}{P_{\mathrm{flux}}}
\newcommand{\epsK}{\varepsilon_K}

%% Integrated paper v10 (2026-11-05). Base: ../研究まとめ_論文草稿_v10.md. Study 1 is restricted to phenomenological
%% noise; its circuit-level code is kept in the repository but not reported here.
%% Blue = to be filled in or revised by the authors (production values, bibliographic data, authors).
%% Numbers: study 1 = stored v7 production data (identical sample and decision rules in v10), 10^4 shots/point;
%% study 2 = v9/v3 production (2.88 M shots); study 3 = 5.54 M shots, v10 re-analysis (results_v10/).
%% Figures: python 論文草稿/統合論文_v10/scripts/make_figures.py (from the stored CSVs). Build: pdflatex main.tex (twice).
\begin{document}

\title{Evaluating Code Switching of the $\mathbb{D}_4$ Clifford-Hierarchy Code\\ with Purpose-Adapted Exact Simulation of Non-Clifford Circuits}

\author{
  \blue{[Authors and affiliations]}
}

\abstract{
The two-dimensional $\mathbb{D}_4$ Clifford-hierarchy stabilizer code of Kobayashi, Zhu and Hsin has transversal logical $T$ and $CS$ gates and, by code switching with the folded surface code, prepares a logical $T$ magic state in $O(d)$ rounds. Evaluating this protocol is a simulation problem: the code's $X$-type stabilizers are dressed by CZ gates, their measurement circuits contain CCZ gates, the switching unitary contains CS and $T^\dagger$ gates, and dressed stabilizers of different colors stop commuting when a flux is present, so that syndromes are not functions of an error configuration. This paper has two contributions. (I) We evaluate code switching at three levels. As a baseline, the finite-size crossings of the charge-sector failure rate of the $\mathbb{D}_4$ memory under phenomenological noise are $2.25\%$, $2.09\%$ and $2.00\%$ ($L=3$--$5$, $5$--$7$, $7$--$9$). For the complete switch FSC $\to\mathbb{D}_4\to T^\dagger\to$ FSC under phenomenological noise the crossings are $2.97\%$ and $2.41\%$ ($d=3$--$5$, $5$--$7$) if the state is idealized just before the switching unitary, but $2.23\%$ and $1.94\%$ otherwise, and at $p=0.8\%$, $d=7$ the residual errors that pass the unitary raise the failure rate more than tenfold relative to the idealized case. Under circuit-level noise, using $5.54\times10^6$ shots, a multinomial treatment of the output infidelity and confidence sequences that remain valid under the stopping rules used, code switching needs at least $1.59$ times the space--time volume of 15-to-1 distillation for an output error of $10^{-4}$ in all ten conditions studied; the single data-certified case requires $2.30$ times. (II) We propose to choose the exact representation of a noisy non-Clifford circuit by the question asked, within the Clifford-stabilizer framework of Takada, Bartlett and Williamson: a code-specific twisted-frame representation (new) for fast memory experiments, an operator-level Pauli-proxy simulation with checkable sufficient conditions for protocols with ideal logical non-Clifford unitaries, and a circuit-level Pauli-proxy simulation for gate-level resource estimates. The three are cross-validated against dense state vectors, brute-force and independent exact engines, and true non-Clifford circuits, with zero discrepancies.
}

\keywords{Code Switching, Clifford-Hierarchy Stabilizer Codes, Magic-State Preparation, Magic-State Distillation, Classical Simulation, Clifford Stabilizers, Just-in-Time Decoding}

\maketitle

%=====================================================================
\section{Introduction}
\label{sec:intro}

Universal fault-tolerant quantum computation needs logical gates outside the Clifford group~\cite{eastin-knill2009}, and for two-dimensional topological Pauli stabilizer codes the constant-depth logical gates are Clifford~\cite{bravyikonig2013}. The standard way to obtain a logical $T$ gate in the surface code is magic-state distillation~\cite{bravyikitaev2005}, whose lattice-surgery layouts and costs have been worked out in detail by Litinski~\cite{litinski2019game,litinski2019magic}. Kobayashi, Zhu and Hsin (KZH) introduced \emph{Clifford-hierarchy stabilizer codes}, whose stabilizer groups contain non-Pauli operators~\cite{kobayashi2025clifford}. In two dimensions, their twisted $\mathbb{Z}_2^3$ gauge theory, equivalent to the non-Abelian $\mathbb{D}_4$ topological order~\cite{iqbal2024nonabelian}, supports transversal logical $T$ and $CS$ gates, and code switching from the folded surface code (FSC) with just-in-time (JIT) flux correction prepares a logical magic state in $O(d)$ rounds~\cite{kobayashi2025clifford,davydova2025knots}. Whether this is cheaper than distillation is an open quantitative question.

Answering it requires simulating noisy circuits that are not Clifford. The $X$-type stabilizers of the $\mathbb{D}_4$ code are Pauli $X$ operators dressed by CZ gates, so their ancilla circuits contain CCZ gates, and the switching unitary contains CS, CS$^\dagger$ and $T^\dagger$. Stabilizer simulators~\cite{gidney2021stim,gottesmanknill1998} cannot run such circuits, and universal simulators~\cite{haenel2026tsim} scale exponentially with the number of non-Clifford gates. More fundamentally, the dressed stabilizers commute only on flux-free states: once an $X$ error creates a flux, dressed stabilizers of different colors can anticommute and randomize each other, and even after the flux is removed a CZ-type residual can randomize them (Sec.~\ref{sec:noncomm}). A previous threshold estimate~\cite{haga2026aqis} replaced the dressed checks by commuting Pauli checks plus independent $Z$ noise and therefore could not represent this. Takada, Bartlett and Williamson showed that circuits of this kind can nevertheless be sampled exactly and efficiently by Clifford-stabilizer simulation~\cite{takada2026exact}.

This paper makes two contributions.

\textbf{(I) An evaluation of code switching for the $\mathbb{D}_4$ code} at three levels, each answering a different question (Sec.~\ref{sec:results}). (a) A memory experiment on the torus gives the finite-size crossings of the dressed-stabilizer (charge-sector) failure rate under phenomenological noise, the baseline for the $\mathbb{D}_4$ phase of the protocol. (b) The complete switch on a planar patch, FSC $\to\mathbb{D}_4\to$ logical $T^\dagger\to$ FSC, under phenomenological noise with ideal unitaries, shows how much the errors that remain at the switching unitary cost. (c) The complete switch as a gate-level circuit under circuit-level noise, compared with 15-to-1 distillation in space--time volume at equal noise, with a statistical analysis that distinguishes what the data certify from what a fitted model supports.

\textbf{(II) Purpose-adapted exact simulation of non-Clifford circuits} (Sec.~\ref{sec:methods}). All three studies are exact samplings within the framework of Ref.~\cite{takada2026exact}, but each uses the representation suited to its question: a \emph{twisted frame} specific to the $\mathbb{D}_4$ code, in which the state is an affine space of flat gauge configurations with a quadratic phase (new; Sec.~\ref{sec:tf}); an \emph{operator-level Pauli proxy}, for which we formulate sufficient conditions of exactness for the code-switching protocol and check them numerically (Sec.~\ref{sec:proxy}); and a \emph{circuit-level Pauli proxy} (Algorithm~2 of Ref.~\cite{takada2026exact}) for the full gate-level circuit, up to $1007$ qubits (Sec.~\ref{sec:circ}). We cross-validate the representations against each other and against non-Clifford references (Table~\ref{tab:validation}) and compare their costs (Table~\ref{tab:methods}).

Section~\ref{sec:background} reviews the code and the protocols, Sec.~\ref{sec:setup} defines noise models, decoders and statistics, Sec.~\ref{sec:discussion} discusses what the results do and do not imply, and Sec.~\ref{sec:conclusion} concludes.

%=====================================================================
\section{Background}
\label{sec:background}

\subsection{The $\mathbb{D}_4$ code}
\label{sec:d4}

We use the triangulated $L\times L$ torus in which every unit square is cut by the diagonal from its lower-left to its upper-right corner: $L^2$ vertices $(x,y)$, $3L^2$ edges $\eh{x,y},\ev{x,y},\ed{x,y}$ oriented right, up and diagonally up, and $2L^2$ triangles. Every edge carries three qubits of colors $c\in\{\rd,\gr,\bl\}$, so $n=9L^2$. A computational-basis state $\ket{a}$, $a\in\GF^n$, is a triple of $\mathbb{Z}_2$ 1-cochains (a \emph{gauge configuration}). The $Z$-type stabilizers are the plaquettes $S^c_{Z;\Delta}=\prod_{e\in\partial\Delta}Z^c_e=(-1)^{(da_c)(\Delta)}$; a configuration is \emph{flat} if all of them are $+1$, and a triangle with $(da_c)(\Delta)=1$ carries a flux of color $c$. The $X$-type stabilizers, one per vertex $v$ and color, are~\cite{kobayashi2025clifford}
\begin{equation}
  S^{\rd}_{X;v}=\Big(\prod_{e\ni v}X^{\rd}_e\Big)(-1)^{\int\tilde v\cup a_\gr\cup a_\bl},
  \label{eq:xstab}
\end{equation}
and similarly for $\gr$ and $\bl$ with $(-1)^{\int a_\rd\cup\tilde v\cup a_\bl}$ and $(-1)^{\int a_\rd\cup a_\gr\cup\tilde v}$, where $\tilde v$ is the indicator 0-cochain of $v$. Evaluated on this lattice, each dressing is a product of two CZ gates on the two colors other than $c$; for example
\begin{equation}
  S^\rd_{X;v}:\ \mathrm{CZ}\big(\gr\,\eh{x,y},\bl\,\ev{x+1,y}\big)\,\mathrm{CZ}\big(\gr\,\ev{x,y},\bl\,\eh{x,y+1}\big)
  \label{eq:czlist}
\end{equation}
for $v=(x,y)$. We write a dressed check $A=(c,v)$ as
\begin{equation}
  S_A=X^{s_A}(-1)^{q_A(a)},\qquad S_A\ket{a}=(-1)^{q_A(a)}\ket{a+s_A},
  \label{eq:generic}
\end{equation}
where $s_A$ marks the six color-$c$ qubits on the edges at $v$ and $q_A(a)=\sum_{(i,j)\in P_A}a_ia_j$ is bilinear over its two CZ pairs $P_A$. Measuring $S_A$ with an ancilla requires six controlled-$X$ gates and two CCZ gates (Fig.~\ref{fig:xstab_circuit}). The planar patch used for code switching (Fig.~4 of Ref.~\cite{kobayashi2025clifford}) has three boundaries; the red and blue layers form an FSC, and the green layer is added during the switch.

\begin{figure}[t]
    \centering
    \input{fig_x_stabilizer.tex}
    \caption{Ancilla circuit for a dressed check $S^{\rd}_{X;v}$: six controlled-$X$ gates for the Pauli part and two CCZ gates for the CZ dressing of Eq.~\eqref{eq:czlist}.}
    \label{fig:xstab_circuit}
\end{figure}

\subsection{Non-commutativity and CZ-type residuals}
\label{sec:noncomm}

Applying Eq.~\eqref{eq:generic} twice gives $S_AS_B=S_BS_A(-1)^{\Delta_{AB}(a)}$ with
\begin{equation}
  \Delta_{AB}(a)=q_A(a+s_B)+q_A(a)+q_B(a+s_A)+q_B(a).
  \label{eq:delta}
\end{equation}
Because the $q$'s are bilinear, $\Delta_{AB}$ is affine; for the cup-product dressing it is a sum of plaquette constraints, $S_AS_B=S_BS_A\prod_{f\in F_{AB}}S^{c_f}_{Z;f}$, non-trivial for exactly three pairs of checks per unit square (27 of the 351 pairs at $L=3$). The checks therefore commute on flat states, as required, but two of them anticommute on a state with a flux on $F_{AB}$. If $S_A$ and $S_B$ anticommute on the current state, then in the sequence $A\to B\to A$ the two outcomes of $A$ agree with probability $\tfrac12$, and this re-randomization repeats as long as the flux survives. A second mechanism exists without flux. For a CZ pair $(q_1,q_2)$ of $S$, $X_{q_1}SX_{q_1}=SZ_{q_2}$. An $X$ error and its correction that together form a product of \emph{bare} $X$ stars, rather than of dressed stabilizers, leave a CZ-type residual (a diagonal Clifford error) that randomizes dressed checks of the other colors although every current generator commutes with every other. Whether the generators currently commute and whether the current state makes a measurement deterministic are different questions. In a diagnostic run with perfect measurements (code-capacity noise, $p=10\%$, $L=5$ and $9$), no flux was ever present before a layer of dressed checks, yet $1.2$--$1.4\%$ of the dressed outcomes were random.

\subsection{Code switching and 15-to-1 distillation}
\label{sec:protocols}

The KZH protocol~\cite{kobayashi2025clifford} is: (1) prepare the FSC in logical $\ket{+}$ and the green layer in $\ket{0}$; (2) measure all $\mathbb{D}_4$ stabilizers for $O(d)$ rounds while correcting fluxes just in time; (3) apply $U=WV$, where $V=\prod\mathrm{CNOT}(\rd\to\gr)\mathrm{CNOT}(\bl\to\gr)$ and $W$ consists of CS, CS$^\dagger$ and boundary $T^\dagger$ gates, which acts as a logical $T^\dagger$; (4) measure the green layer in the $Z$ basis, apply the CZ part of the green $X$ stabilizers conditioned on the outcomes, and continue with $O(d)$ FSC rounds. The output is $T^\dagger\ket{+}_L$. Qubits that are fixed by a single $Z$ at the boundary are removed as ideally fixed (reduced model); the lattices of Secs.~\ref{sec:res2} and \ref{sec:res3} agree generator by generator for $d=2,3,5,7,9$ after this removal. The 15-to-1 protocol~\cite{bravyikitaev2005} consumes 15 noisy $T$ states and outputs one with error $\approx35\epsilon^3$; we use the 11-tile, $11\tau$ layout of Ref.~\cite{litinski2019game} and the $(15\text{-to-}1)_{d_X,d_Z,d_m}$ and two-level protocols with the error-budget model of Ref.~\cite{litinski2019magic}.

%=====================================================================
\section{Purpose-adapted exact simulation}
\label{sec:methods}

\subsection{Common principle and choice of representation}

If a Pauli error $E_P$ is conjugated through a Clifford factor $C$, $E_PCE_P^\dagger=D_C(E_P)C$, the \emph{dressing} $D_C(E_P)$ is a Pauli operator for the circuits considered here. When every dressed operator commutes or anticommutes with the current stabilizers, each measurement is deterministic or uniformly random, and the update of the Clifford-stabilizer group reduces to that of a Pauli stabilizer group~\cite{takada2026exact}. The sampling is then exact. This is a statement about the simulation, independent of the optimality of the decoder, the fault tolerance of the protocol or the existence of an asymptotic threshold.

Within this framework there is more than one way to hold the state, and the best one depends on the question. Our proposal is to pick the representation from three properties of the task: which non-Clifford elements occur (only in stabilizers, also in ideal logical unitaries, or as noisy gates), which output is needed (failure rates at many sizes, a logical channel class, or an output infidelity at low $p$), and which reference can validate it. Table~\ref{tab:methods} summarizes the three representations used in this work.

\begin{table*}[t]
\centering
\caption{The three exact representations, the questions they answer, and their measured costs. Times per shot on one CPU core, including decoding. Twisted frame and operator-level proxy for the memory experiment: phenomenological noise at $p=2\%$, measured for this paper on one core of an Intel Xeon (2.1\,GHz) (\texttt{scripts/bench\_study1\_representations.py}); ``same driver'' runs both through the same per-shot Python protocol, ``batched'' is the production engine. Code switching: medians of the stored production runs. \blue{[Confirm the hardware of the stored runs of studies 2 and 3.]}}
\label{tab:methods}
\small
\renewcommand{\arraystretch}{1.15}
\begin{tabular}{p{0.15\linewidth}p{0.25\linewidth}p{0.25\linewidth}p{0.25\linewidth}}
\toprule
 & twisted frame (Sec.~\ref{sec:tf}) & operator-level Pauli proxy (Sec.~\ref{sec:proxy}) & circuit-level Pauli proxy (Sec.~\ref{sec:circ}) \\
\midrule
question & $\mathbb{D}_4$ memory: crossings, many sizes and $p$ & switch with ideal $U$: logical class, cost of the residual at $U$ & gate-level switch: output infidelity, resource comparison \\
non-Clifford content & dressed stabilizers (CZ factors), measured as projections & dressed stabilizers; $U$ with CS, CS$^\dagger$, $T^\dagger$ (ideal) & CCZ in check circuits; CS, CS$^\dagger$, $T^\dagger$ in $U$; all gates noisy \\
state & affine space $B$ of flat configurations, quadratic form $Q$, $X$ frame $x$ (all $\GF$) & Stim tableau of the proxy state, error frame $E=X^x\mathrm{diag}(\omega^\gamma)$ & Clifford circuit in Stim, forward-propagated $X$ components, inserted dressings \\
scope & $\mathbb{D}_4$ code only (needs the flat-subspace structure) & protocols satisfying C1--C3 & circuits satisfying the applicability conditions \\
output per shot & outcome probabilities, exact logical fidelity $F$ & logical class $\{I,X,Y,Z,S,\text{collapsed}\}$ & infidelity $\in\{0,\tfrac12,1\}$ \\
time per shot & batched: $1.0$, $1.8$, $4.6$\,ms ($L=5,7,9$); same driver: $4.2$, $5.1$, $11.9$, $24.6$\,ms ($L=3,5,7,9$) & memory, same driver: $3.7$, $17.6$, $63.5$, $208$\,ms ($L=3,5,7,9$); switch: $1.8$, $9.1$, $21$\,ms ($d=3,5,7$) & scenario B: $2.5$, $9.0$, $31$, $108$\,ms; A: $3.3$, $17$, $54$, $134$\,ms ($d=3,5,7,9$) \\
largest size & $L=9$ ($729$ qubits); $L\ge11$ feasible & $d=7$ & $d=9$ ($1007$ qubits) \\
\bottomrule
\end{tabular}
\end{table*}

\subsection{Twisted frame}
\label{sec:tf}

On the torus, a flat configuration of one color with vertical holonomy fixed to zero is $a_c=d\mu_c+h_c\eta$, with $\mu_c$ a 0-cochain ($\mu_c(0,0)=0$), $h_c\in\GF$ the horizontal holonomy and $\eta$ the indicator of the edges crossing a fixed vertical line. These configurations form the space $C_0=\{Mt:\ t\in\GF^k\}$, $k=3L^2$. Let $\ket{\psi_0}=\sum_{b\in C_0}(-1)^{\phi_0(b)}\ket{b}$ be a code state. Because every dressed check stabilizes it,
\begin{equation}
  \phi_0(b+s)+\phi_0(b)=q(b+s)\qquad(b\in C_0);
  \label{eq:stabcond}
\end{equation}
the cubic phase $\phi_0$ itself is never needed.

\begin{proposition}
Let $\ket{\psi}$ be obtained from $\ket{\psi_0}$ by Pauli operators, CZ gates, plaquette measurements and projective measurements of dressed checks (possibly times a Pauli $Z$ on a dressing qubit). Then, up to normalization,
\begin{equation}
  \psi(a)=(-1)^{\phi_0(a+x)+Q(a+x)}\,\mathbf{1}\big[a+x\in B\big],
  \label{eq:tf}
\end{equation}
where $x\in\GF^n$ is the accumulated $X$ frame, $B\subseteq C_0$ is an affine subspace and $Q$ is a quadratic form over $\GF$. Every outcome has probability $0$, $\tfrac12$ or $1$.
\end{proposition}

\begin{proof}[Proof sketch]
$X$ errors change $x$; $Z$ errors and CZ gates add linear or quadratic terms to $Q$; plaquette outcomes are $(dx)(f)$ because $B$ is flat. For a dressed check with shift $s$ and $b=a+x$, at points with $b,b+s\in B$ Eq.~\eqref{eq:stabcond} gives the amplitude ratio $(S\psi)(a)/\psi(a)=(-1)^{f(b)}$ with the affine function $f(b)=q(b+s+x)+q(b+s)+Q(b+s)+Q(b)$. (i) If $s$ lies in the direction space of $B$, the outcome is deterministic when $f$ is constant on $B$, and otherwise uniform with $B\to B\cap\{f=m\}$. (ii) Otherwise $B$ and $B+s$ are disjoint, $\langle S\rangle=0$, the outcome is uniform, $B\to B\cup(B+s)$, and $Q\to Q+\lambda(b)(f(b)+m)$, where $\lambda$ is the affine indicator of $B+s$. Both preserve Eq.~\eqref{eq:tf}.
\end{proof}

Case (ii) is the non-commutativity of Sec.~\ref{sec:noncomm}, and the $f$ of case (i) contains the CZ-type residual through the frame $x$. Storing $x$, a row-reduced basis of the constraints of $B$ and $Q$ as bit-packed $k\times k$ matrices, a measurement costs $O(k^2)$ bit operations. At the end, the phases of the eight holonomy sectors give the exact logical fidelity. The representation is a special case of the Clifford-stabilizer simulation~\cite{takada2026exact}: the group commutator of two dressed checks is a product of plaquettes, so every measurement is effectively commuting or anticommuting. What is new is that the state is held directly as an affine space with a quadratic phase. This is fast, because no tableau is propagated, and the representation yields outcome probabilities and fidelities as exact numbers rather than samples. Its limitation is that it uses the flat-subspace structure of this code.

\subsection{Operator-level Pauli proxy}
\label{sec:proxy}

For the switch with ideal $U$ we keep the accumulated errors and corrections as $E=X^x\mathrm{diag}(\omega^{\gamma(a)})$ with $\gamma(a)=\sum_il_ia_i+4\sum_{i<j}Q_{ij}a_ia_j \bmod 8$ and $\omega=e^{i\pi/4}$, and measure a dressed generator $M=X^s\Delta$ as $E^\dagger ME$ on a Pauli stabilizer state held by Stim. $E^\dagger ME$ is $M$ times a Pauli dressing; replacing the CZ factor $\Delta$ by its Pauli proxy, the identity, leaves a Pauli operator. Writing the physical state as $E\,C\ket{\mathrm{frame}}$ with $C=\mathrm{diag}((-1)^{\phi_0})$, the reduction is exact if
\begin{enumerate}[leftmargin=*,itemsep=0pt,label=(C\arabic*)]
\item $CX^{s_G}C^\dagger=G$ on the flat subspace for every dressed generator $G$;
\item $C$ acts trivially on FSC states;
\item $UC=CVT_h$ on the flat subspace, with $T_h=\mathrm{diag}(\omega^{7h})$ the logical $T^\dagger$ kept as bookkeeping (a pointwise form of Theorem~2 of Ref.~\cite{kobayashi2025clifford}).
\end{enumerate}
C1 follows from the definition of $\phi_0$ and C3 from that theorem; the code checks all three on random flat configurations at every $d$ used (an implementation check, not a proof). The proxy of $U$ is then $V$, a Clifford circuit. The switch-out applies the branch-flipping operator $S^g_X(\lambda)$ of the green measurement to $E$ (Algorithm~3 of Ref.~\cite{takada2026exact}), and the logical result is classified by the expectation values of the proxies of the logical Clifford stabilizers of the target state (Sec.~VI\,D of Ref.~\cite{takada2026exact}). The same representation also runs the memory experiment, with the initial frame the uniform superposition of flat configurations; we use it there only for cross-validation.

\subsection{Circuit-level Pauli proxy}
\label{sec:circ}

For the gate-level circuit we sample a Pauli fault configuration per shot and run in Stim the Clifford part of the true circuit, i.e.\ the circuit with the diagonal third-level gates $D$ (CCZ, CS, CS$^\dagger$, $T^\dagger$) removed. The $X$ components $e$ of the faults and of the applied corrections are propagated forward; where $e$ meets the support of a removed $D$, the dressing $C_e=D^\dagger X^eDX^e$ is inserted in its place (Algorithm~2 of Ref.~\cite{takada2026exact}). We check the applicability conditions for $L=2,\dots,7,9$: the $X$ and CZ factors of each generator have disjoint supports, dressings are single-qubit $Z$, group commutators lie in the span of the $Z$ stabilizers, primary-support exactness holds, $U$ and the corrections contain only CNOTs and diagonal third-level gates, and the logical stabilizers of the output are Clifford.

\subsection{Cross-validation}

Because the reduction is exact, any disagreement with a reference that does not use it is an implementation error. Table~\ref{tab:validation} lists the checks. The memory experiment is compared with a dense state vector over the flat configurations ($2^{25}$ or $2^{27}$ amplitudes, computed with JAX~\cite{jax2018github}), and the twisted frame and the operator-level proxy, which share no state representation, are driven in lock step up to $L=9$. Deliberately removing the Pauli dressing, or the CZ part of the error frame, makes this comparison fail within the first trajectories. The two engines do share the geometry, the protocol and the decoders. The switch with ideal $U$ is compared with three references (brute-force state vectors, an independent exact engine of the twisted-frame family, and a replay of the production path). The circuit-level engine is compared with the true non-Clifford circuit at $L=1,2$.

\begin{table}[t]
\centering
\caption{Cross-validation of the exact representations. ``Outputs'' counts every compared quantity (outcome and its determinism, $Z$ checks, $X$ frame, final class).}
\label{tab:validation}
\small
\renewcommand{\arraystretch}{1.1}
\begin{tabular}{p{0.36\linewidth}p{0.55\linewidth}}
\toprule
comparison & result \\
\midrule
twisted frame vs dense state vector, memory, $L=3$, phenomenological, $p=4\%$ & 40 trajectories, 216 random outcomes; max.\ deviation of probabilities $1.9\times10^{-9}$, of fidelities $1.2\times10^{-7}$; independent samples: $\sum z^2=6.49$ (6 d.o.f., $p=0.37$) \\
twisted frame vs operator-level proxy, memory, $L=3$--$9$ & phenomenological: 3\,600 trajectories, $9.5\times10^4$ random outcomes; all noise models of the code: 14\,400 trajectories, $5.3\times10^6$ measurements; identical probabilities and fidelities \\
operator-level proxy vs brute-force state vector, switch, $d=2$ & 5\,400 shots, 683\,004 outputs, 0 mismatches \\
vs independent exact engine, switch, $d=3,4$ & 7\,200 shots, 4\,069\,644 outputs, 0 mismatches \\
vs replay of the production path, $d=2$ & 1\,800 shots, 227\,640 outputs, 0 mismatches \\
C1--C3, flat commutation, $p=0$ & 0 violations ($d=3,5,7$); 0 failures at $p=0$ \\
circuit-level proxy vs true non-Clifford circuit & $L=1,2$ (incl.\ all adaptive paths at $L=2$); 189 checks PASS, 0 FAIL \\
\bottomrule
\end{tabular}
\end{table}

\paragraph{Costs.}
Table~\ref{tab:methods} shows that the choice matters in practice. Run through the same per-shot driver, the operator-level proxy is $0.9$, $3.5$, $5.3$ and $8.5$ times slower than the twisted frame at $L=3,5,7,9$, and the batched twisted-frame engine needs $4.6$\,ms per shot at $L=9$, $45$ times less than the proxy. The $10^5$-shot production of the memory experiment (about $10^7$ shots including the control variants) therefore needs about half a day of CPU time per noise model instead of weeks. Conversely, the twisted frame cannot represent the switching unitary and the gate-level circuit, for which the two proxy representations remain practical up to $d=9$.

%=====================================================================
\section{Protocols, noise, decoders and statistics}
\label{sec:setup}

\paragraph{Noise models.}
$\mathcal D_k(p)$ is the $k$-qubit depolarizing channel (each non-identity Pauli with probability $p/(4^k-1)$). Table~\ref{tab:noise} gives the models. Studies 1 and 2 use phenomenological noise with ideal gates; study 3 uses the circuit-level model of App.~A2 of Ref.~\cite{takada2026exact}, in which the CCZ gate is a native gate and faults inside a $T$-gate decomposition are not considered. JIT corrections and feed-forward Pauli corrections are noiseless.

\begin{table}[t]
\centering
\caption{Noise models. S0, S1, S4: FSC preparation, $\mathbb{D}_4$ phase, FSC phase (Sec.~\ref{sec:res3}).}
\label{tab:noise}
\small
\renewcommand{\arraystretch}{1.1}
\begin{tabular}{p{0.22\linewidth}p{0.69\linewidth}}
\toprule
study 1, memory (phenomenological) & $\mathcal D_1(p)$ on all data qubits per round; every plaquette and dressed-check outcome flipped with probability $p$; $R=L$ rounds \\
study 2, switch (phenomenological) & $\mathcal D_1(p)$ per switch-in round, once after $U$, and on red/blue per FSC round; all check outcomes and the green $Z$ readout flipped with $p$; $d$ switch-in rounds, $d$ FSC rounds and one perfect round; $U$ and corrections ideal \\
study 3, switch, circuit level B & $\mathcal D_k(p)$ after every $k$-qubit gate (CX: 15, CCZ: 63 Paulis); preparation and measurement flips $p$; no idle noise; $d$ rounds in each of S0, S1, S4 \\
study 3, circuit level A & as B, plus idle $\mathcal D_1(p)$ per time step \\
\bottomrule
\end{tabular}
\end{table}

\paragraph{Decoders.}
In all studies new fluxes are paired in real time by minimum-weight perfect matching (PyMatching~\cite{higgott2025pymatching}) on a space--time graph that is open to the future, and the spatial part of every decided correction is applied at once~\cite{bombin2018jit,brown2020jit,kobayashi2025clifford}. An event left unpaired costs a waiting weight per round, so an isolated flipped plaquette waits for its partner in the next round, whereas two fresh adjacent events are corrected before any dressed check sees them. This is the JIT idea of Refs.~\cite{kobayashi2025clifford,davydova2025knots}, not the decoder of Ref.~\cite{davydova2025knots}. Study 1 corrects only pairs of real events. Study 2 uses the planar graph with spatial boundaries and also corrects the spatial part of paths to the future boundary. Study 3 re-decodes the full history every round and joins the switch-out $Z^g$ readout to the green fluxes. Across $U$, study 3 uses one space--time graph for the whole protocol. In study 2 without idealization, fluxes pending at $U$ are decoded with the switch-out readout (green) and the FSC rounds (red, blue), and red/blue charges with the switch-in and FSC rounds, which is possible because $Z^r$ and $Z^b$ commute with $U$ and with $S^g_X$. Green charges are cleaned up before $U$, and their remainder becomes pairs of red/blue charges under $V$; correlations and hyperedges are not modeled. Charges in study 1 are decoded by space--time MWPM with the detector error model of a bare-check Stim circuit as prior; study 3 decodes on the error model of the proxy circuit.

\paragraph{Failure metrics.}
Study 1 reports $\Pch$, the charge-sector failure probability conditioned on flux success, together with $\Pfl$ and $P_{\mathrm{comb}}=\Pfl+(1-\Pfl)\Pch$. The conditioning uses the true error frame of the simulation and is an analysis quantity, not an experimental post-selection. The initial state is the superposition of the eight horizontal holonomy sectors, and the shot fails with probability $1-F$. Study 2 reports $P_L=1-P(I)$, the fraction of shots whose final logical class is not the target, which differs from the mean infidelity (the classes $X$, $Y$, $S$ and collapsed contribute $\tfrac12$ to the latter). Study 3 reports the output infidelity $\epsK$ of $T^\dagger\ket{+}_L$.

\paragraph{Statistics.}
Failure rates in studies 1 and 2 carry binomial 95\% Agresti--Coull intervals~\cite{agresti1998}. Crossings are obtained by log--log interpolation between neighboring sizes; their ranges are the 2.5--97.5\% quantiles over failure counts redrawn from Jeffreys posteriors, which describe the sampling variation of the crossing position, not an interval for an asymptotic threshold. In study 3 the per-shot infidelity is $X\in\{0,\tfrac12,1\}$ and the counts $(n_0,n_{1/2},n_1)$ are multinomial, $\hat\varepsilon=(n_1+n_{1/2}/2)/N$ and $\mathrm{Var}(\hat\varepsilon)=(p_1+p_{1/2}/4-\varepsilon^2)/N$. Treating $2S=2n_1+n_{1/2}$ as Poisson underestimates the variance by about a factor of two when $n_1$ dominates. Since the Monte Carlo runs stopped at failure-count and time limits and points were extended later, fixed-$N$ intervals do not apply. We use the Beta$(\tfrac12,\tfrac12)$-mixture confidence sequence for a mean of a $[0,1]$-valued loss~\cite{howard2021cs},
\begin{equation}
  E_N(\mu)=\frac{B(S+\tfrac12,N-S+\tfrac12)}{B(\tfrac12,\tfrac12)\,\mu^S(1-\mu)^{N-S}},\quad \{\mu:\ E_N(\mu)<M/\alpha\},
  \label{eq:cs}
\end{equation}
with $S=\sum X$. $E_N(\mu)$ is a non-negative supermartingale under the null, so by Ville's inequality~\cite{ville1939} the interval holds at any stopping time. The $M=4$ simulated candidate sizes of one decision (scenario, $p$) are covered simultaneously by a Bonferroni correction, so that the reported size and the statement that smaller sizes fail hold together, for all targets of that row. A size is \emph{data-certified} if the upper end is below the target and \emph{data-rejected} if the lower end is above it. A stricter, table-wide family ($M=200$) is reported as a sensitivity. For extrapolation we fit a multinomial log-linear model with a bulk and a switching-window component and call a size \emph{model-supported} if at least 95\% of 200 parametric-bootstrap fits reach the target. Model support is not a certificate and never overrides a data decision.

%=====================================================================
\section{Results: evaluation of code switching}
\label{sec:results}

\subsection{Baseline: $\mathbb{D}_4$ memory under phenomenological noise}
\label{sec:res1}

Each of the $R=L$ rounds consists of data noise, plaquette measurements, JIT flux correction and the dressed checks measured color by color ($\rd\to\gr\to\bl$). The run starts from the noiseless code state and ends with a perfect plaquette round, closing of the remaining fluxes, removal of a trivial $X$ frame by applying dressed stabilizers as unitaries, and a perfect dressed round. Figure~\ref{fig:memory} and Table~\ref{tab:cross1} show the results. The crossings move to lower $p$ with $L$, by $0.16$ and $0.09$ percentage points, so they are finite-size estimates; asymptotic thresholds are expected to be somewhat lower. Near the crossing $\Pfl$ is more than an order of magnitude smaller than $\Pch$ (at $p\approx2.0\%$, $L=9$: $0.33\%$ versus $12.8\%$), so the crossing is not an artifact of conditioning on flux success. With measurement errors, fluxes survive into a layer of dressed checks in $66\%$ of the rounds ($L=7$, $p=2\%$), and $1.9$--$2.7\%$ of all dressed outcomes near the crossing are random. Only 2 of the $4\times10^5$ shots had a final fidelity strictly between 0 and 1.

\begin{figure}[t]
  \centering
  \includegraphics[width=\linewidth]{figures/fig_memory.pdf}
  \caption{$\mathbb{D}_4$ memory, phenomenological noise: $\Pch$ (solid) and $\Pfl$ (dashed) versus $p$ for $L=3,5,7,9$, $10^4$ shots per point, 95\% Agresti--Coull intervals. \blue{[Replace by the $10^5$-shot production of \texttt{D4\_v10\_main.ipynb} (phenomenological); regenerate with \texttt{scripts/make\_figures.py} after pointing it to the new CSV.]}}
  \label{fig:memory}
\end{figure}

\begin{table}[t]
\centering
\caption{Crossings of $\Pch$ (\%) under phenomenological noise, $10^4$ shots per point; in parentheses the 2.5--97.5\% range over resampled data. \blue{[Replace by the $10^5$-shot production.]}}
\label{tab:cross1}
\small
\begin{tabular}{lccc}
\toprule
 & $L=3$--$5$ & $5$--$7$ & $7$--$9$\\
\midrule
crossing & \blue{2.247} & \blue{2.088} & \blue{1.998}\\
range & \blue{2.174--2.302} & \blue{2.014--2.133} & \blue{1.943--2.065}\\
\bottomrule
\end{tabular}
\end{table}

\paragraph{Control experiments.}
With identical noise, protocol and decoders we ran two variants \blue{[pilot runs, $2\times10^4$ shots per point; replace by the production runs]}: \emph{untwisted}, the dressing removed (three independent toric codes), and \emph{$\kappa$-injection}, untwisted plus the independent pair-$Z$ noise of the previous surrogate~\cite{haga2026aqis}. The untwisted crossings are \blue{$3.11$, $2.88$, $2.70$ and $2.72\%$} ($L=3$--$5$ to $9$--$11$), about \blue{$25$--$30\%$} above the $\mathbb{D}_4$ values (\blue{$2.25$, $2.12$, $1.96$, $1.89\%$} in the same pilot). This is the total difference between the chosen dressed and undressed models, not an isolated effect of non-commutativity. The $\kappa$-injection crossings (\blue{$2.31$, $2.23\%$}) happen to lie near the $\mathbb{D}_4$ ones, but its failure rates are higher at every point (\blue{$0.20$, $0.15$, $0.12$ versus $0.14$, $0.10$, $0.08$} at $p=1.8\%$, $L=3,5,7$), so independent noise does not reproduce the $\mathbb{D}_4$ curves. The previous estimate of $2.4\%$~\cite{haga2026aqis} differs from ours by the exact dressing, the flux correction, and changes of the initial state, termination and metric together; the difference cannot be attributed to one of them.

\subsection{Code switching under phenomenological noise}
\label{sec:res2}

We use the triangular patch with $(d+1)\times(d+1)$ vertices, where $d$ counts the squares along a side; no claim is made about the fault distance of the whole protocol. A shot runs $d$ switch-in rounds, the boundary handling, the clean-up of charges, $U$ (case \cs{in\_T\_out}), the switch-out (one noise layer, a single green $Z$ readout, green correction), $d$ FSC rounds and one perfect round, and ends with the logical class. Two further cases isolate parts of the protocol: \cs{in\_only} reads out the $\mathbb{D}_4$ code ideally, and \cs{in\_out} switches back without $U$. The \emph{ideal} boundary inserts a noiseless complete round just before $U$, an optimistic reference that decouples switch-in and switch-out; the \emph{noisy} boundary does not, so residual $X$ errors pass $W$ and become non-Pauli, and green charges spread under $V$. The production run covers $d=3,5,7$ and eight values of $p$ between $0.8\%$ and $5\%$, with $2\times10^4$ shots per point, i.e.\ six conditions and $2.88\times10^6$ shots (Fig.~\ref{fig:switching}, Table~\ref{tab:switch}).

\begin{figure*}[t]
  \centering
  \includegraphics[width=\linewidth]{figures/fig_switching.pdf}
  \caption{Code switching under phenomenological noise: $P_L$ versus $p$ for $d=3,5,7$ with the ideal (solid, filled) and noisy (dashed, open) boundary before $U$. (a) Switch-in and ideal $\mathbb{D}_4$ readout; (b) switch-in and switch-out without $U$; (c) the full switch with logical $T^\dagger$. $2\times10^4$ shots per point, 95\% intervals.}
  \label{fig:switching}
\end{figure*}

\begin{table*}[t]
\centering
\caption{Code switching, phenomenological noise: crossings (\%, 2.5--97.5\% range) and $P_L$ at $p=0.8\%$.}
\label{tab:switch}
\small
\begin{tabular}{llcc|ccc}
\toprule
 & & \multicolumn{2}{c|}{crossing} & \multicolumn{3}{c}{$P_L$ at $p=0.8\%$}\\
case & bound. & $d=3$--$5$ & $5$--$7$ & $d=3$ & $5$ & $7$\\
\midrule
\cs{in\_T\_out} & ideal & 2.97 (2.89--3.10) & 2.41 (2.25--2.54) & 0.0111 & 0.0015 & 0.0004\\
 & noisy & 2.23 (2.12--2.35) & 1.94 (1.83--2.01) & 0.0213 & 0.0078 & 0.0051\\
\cs{in\_out} & ideal & 3.13 (3.02--3.21) & 2.24 (2.16--2.37) & 0.0113 & 0.0016 & 0.0003\\
 & noisy & 2.81 (2.74--2.89) & 2.23 (2.15--2.36) & 0.0121 & 0.0021 & 0.0005\\
\cs{in\_only} & ideal & 2.33 (2.21--2.43) & 1.91 (1.73--2.01) & 0.0047 & 0.0009 & 0.0006\\
 & noisy & 2.10 (1.99--2.24) & 1.74 (1.52--1.88) & 0.0066 & 0.0020 & 0.0016\\
\bottomrule
\end{tabular}
\end{table*}

Without $U$ (\cs{in\_out}) the noisy and ideal boundaries differ little. With $U$ they separate. At $p=0.8\%$ and $d=7$ the noisy failure rate is $12.6$ times the ideal one, and going from $d=5$ to $7$ improves it by $1.5$ times instead of $3.8$ times. Summed over all points, the noisy boundary produces about five times more $X$ and $Y$ classes ($18\,291$ and $16\,691$ versus $3\,823$ and $2\,960$) and $1\,527$ collapsed outcomes, which never occur with the ideal boundary. Under the present protocol and decoder, errors left at $U$ are thus amplified by $W$. Because the decoder still approximates the green charges, this is not the intrinsic cost of $U$ under optimal decoding. All crossings decrease with $d$. The \cs{in\_only} rates are not a lower bound for the full protocol: they read out the $\mathbb{D}_4$ code itself and include its flux-sector failures ($20\,732$ of the $4.8\times10^5$ shots for either boundary), a class that does not occur in the cases that switch back to the FSC.

\subsection{Circuit-level resources versus 15-to-1 distillation}
\label{sec:res3}

\paragraph{Circuit and cost.}
The gate-level protocol consists of S0 (FSC preparation, $d$ rounds), S1 (green in $\ket0$, $d$ rounds of all $\mathbb{D}_4$ stabilizers, followed by two $Z$-only rounds), $U$, the switch-out (green readout, decoding, conditional CZ layer), S4 ($d$ FSC rounds) and a noiseless evaluation round. A dressed check uses one ancilla, six CNOTs and two CCZ gates, 28 time steps per $X$-type round. Every gate, preparation and measurement takes one step, and a surface-code round takes six. The circuit has $12d^2+4d-1$ qubits, $62d+21$ steps and space--time volume $632d^3+454d^2+8d-14$, which agree with counts from the circuit for $d=3,\dots,11$. Classical latency, conversion from the FSC to a standard patch, and state injection are not counted. On the distillation side, we take the minimum volume over 16\,320 candidate factories: the layout of Ref.~\cite{litinski2019game} and the protocols of Ref.~\cite{litinski2019magic} with its error-budget model (reimplemented from its public code; eight published values are reproduced), and rotated-surface-code logical error rates calibrated under the same noise (A: $A=0.079$, $p_{\mathrm{th}}=0.0085$; B: $A=0.040$, $p_{\mathrm{th}}=0.0119$). These are central predictions without the statistical guarantee of the Kobayashi side.

\paragraph{Single faults.}
A scan of all single faults, weighted by probability, gives the first-order coefficient $c_1$ ($\epsK\approx c_1p$): $15.7$ (A) and $4.16$ (B) at $d=3$, $0.88$ (A) and $0.27$ (B) at $d=5$ within the switching window, and $0$ for B at $d=7$. This circuit with this decoder therefore does not correct every single fault at $d=5$. The failures concentrate in the window: in B, the CCZ gadgets of the last $X$-type round of the $\mathbb{D}_4$ phase, and in A, the second $Z$ round just before $U$. With only one $Z$ round before $U$, $c_1=3.1$ (B, $d=5$), so the second round reduces the single-fault contribution about tenfold. The opposite holds under phenomenological noise: in a pilot of study 2, an additional pre-$U$ round doubled to tripled the failure rate, because there every round adds data noise to all qubits. The window must be designed together with the noise model.

\paragraph{Output infidelity.}
Figure~\ref{fig:epsK} shows $\epsK$ from $5.54\times10^6$ shots. In B it drops by more than an order of magnitude from $d=3$ to $5$, but improves little for $d\ge5$. Diagnostic runs with a noiseless window differ from these by an approximately constant amount at $d=7,9$, which shows that errors from two or more faults in the switching window remain. With idle noise (A), $\epsK$ is several times larger and does not decrease with $d$ at $p=10^{-3}$.

\begin{figure*}[t]
  \centering
  \includegraphics[width=\linewidth]{figures/fig_epsK.pdf}
  \caption{Output infidelity of circuit-level code switching versus $d$ for five values of $p$ (points slightly offset horizontally). Bars: data intervals of the decision family [Eq.~\eqref{eq:cs}, $M=4$, 95\%]; tri-markers: upper ends for points without failures. Dotted: target $10^{-4}$.}
  \label{fig:epsK}
\end{figure*}

\paragraph{Comparison.}
Table~\ref{tab:cost} lists the comparison at target $10^{-4}$. In each row, the smallest size not rejected by the data, $d_{\mathrm{low}}$, gives a lower bound on the Kobayashi volume, because the volume increases with $d$. Unsimulated sizes count as not rejected.
\begin{itemize}[leftmargin=*,itemsep=0pt]
\item Only one condition is data-certified: B, $p=10^{-4}$, $d=5$ ($N=414\,400$, $S=12$, upper end $8.8\times10^{-5}$). Since $d=3$ is data-rejected, the minimum is resolved too, and code switching needs \textbf{2.30 times} the volume of the 15-to-1 factory.
\item In all ten conditions at target $10^{-4}$, the data lower bound is at least \textbf{1.59 times} the distillation volume (smallest for B, $p=3\times10^{-4}$).
\item The model supports $d=7$ for B at $p=10^{-4}$ and $2\times10^{-4}$, i.e.\ $6.1$ and $4.2$ times the distillation volume.
\item At target $10^{-6}$ every simulated size $d\le9$ is data-rejected except for B, $p=10^{-4}$, where $d_{\mathrm{low}}=7$. The data lower bounds are $1.9$--$10.2$ times the distillation volume in nine of the ten conditions; in the remaining one (A, $p=10^{-3}$) the lower bound, fixed by $d_{\mathrm{low}}=11$, is $0.90$ times. Because $d\ge11$ was not simulated, these bounds are weak. For targets $10^{-8}$ and $10^{-10}$ the bound falls below one in 3 and 7 of the 10 conditions, respectively, where the data do not decide. The model supports no size $d\le61$ for any target $\le10^{-6}$, because the window component does not decrease with $d$.
\item Under the table-wide family ($M=200$) all 40 conditions are unresolved.
\end{itemize}
\blue{[Optional: after the fixed-shot extensions of B, $p=2\times10^{-4}$, $d=7$ and A, $p=10^{-4}$, $d=7$ (and $d\ge11$ for targets $\le10^{-6}$), update Fig.~\ref{fig:epsK}, Table~\ref{tab:cost} and this list with \texttt{DATA\_MODE='production'}.]}

\begin{table}[t]
\centering
\caption{Space--time volume at target $\epsK\le10^{-4}$. $V_{\mathrm{L}}$: minimum 15-to-1 volume (qubit$\cdot$steps, central model). Data: certified $d$ and ratio, or the lower bound from $d_{\mathrm{low}}$. Model: smallest model-supported $d$ (ratio).}
\label{tab:cost}
\small
\setlength{\tabcolsep}{3pt}
\begin{tabular}{ccrll}
\toprule
sc. & $p$ & $V_{\mathrm{L}}$ & data & model\\
\midrule
B & $10^{-4}$ & $3.93\times10^4$ & $d=5$: $2.30$ (certified) & $d=7$: $6.09$\\
B & $2\times10^{-4}$ & $5.66\times10^4$ & $d\ge5$: $\ge1.60$ & $d=7$: $4.22$\\
B & $3\times10^{-4}$ & $5.68\times10^4$ & $d\ge5$: $\ge1.59$ & --\\
B & $5\times10^{-4}$ & $8.92\times10^4$ & $d\ge7$: $\ge2.68$ & --\\
B & $10^{-3}$ & $9.37\times10^4$ & $d\ge11$: $\ge9.57$ & --\\
A & $10^{-4}$ & $5.65\times10^4$ & $d\ge7$: $\ge4.23$ & --\\
A & $2\times10^{-4}$ & $5.69\times10^4$ & $d\ge9$: $\ge8.75$ & --\\
A & $3\times10^{-4}$ & $8.96\times10^4$ & $d\ge11$: $\ge10.0$ & --\\
A & $5\times10^{-4}$ & $9.31\times10^4$ & $d\ge11$: $\ge9.63$ & --\\
A & $10^{-3}$ & $3.92\times10^5$ & $d\ge11$: $\ge2.29$ & --\\
\bottomrule
\end{tabular}
\end{table}

Within the circuit, schedule, decoder and cost model constructed here, 15-to-1 distillation therefore reaches the target with a smaller volume in every condition the data can decide, and the model does not support code switching reaching $10^{-6}$ or below. This is not a statement about an optimized Kobayashi protocol.

%=====================================================================
\section{Discussion}
\label{sec:discussion}

\paragraph{What limits code switching here.}
The three levels of the evaluation point to the same place: the switching window. Under phenomenological noise, the errors left at $U$ cost more than the $\mathbb{D}_4$ phase itself, and the gap to the idealized boundary widens with $d$. Under circuit-level noise, a single fault in the window is not always corrected at $d=5$, and the two-fault contribution in the window sets a floor that larger $d$ does not remove. The other two costs are structural: an $X$-type round of the $\mathbb{D}_4$ phase takes 28 steps, which gives the $632d^3$ volume, and idle noise accumulates during it (scenario A). Changes that could alter the conclusion include decoding across the window (the JIT decoder of Ref.~\cite{davydova2025knots}, hyperedge-aware or flux-history-aware decoders~\cite{magdalena2026}), flagged CCZ gadgets, the actual CCZ error rate relative to two-qubit gates, costs not counted here, and a circuit-level simulation of the distillation side.

\paragraph{Comparing the numbers.}
The memory crossing ($\approx2.0$--$2.25\%$) is a crossing of $\Pch$ on the torus, conditioned on flux success. The switch crossings are crossings of $P_L$ for the complete protocol on a patch with boundaries, and $\epsK$ is an infidelity at fixed $(d,p)$. The memory and \cs{in\_only} crossings are close ($1.9$--$2.3\%$), but these are not two measurements of one threshold.

\paragraph{What the method choice enabled.}
The twisted frame made the memory sweep cheap enough for $10^5$ shots per point and $L=9$; the operator-level proxy made $U$ representable with checkable conditions; the circuit-level proxy reached $1007$ qubits with every gate noisy. Their agreement on the overlap (memory, $L\le9$) and with the non-Clifford references (Table~\ref{tab:validation}) is what lets the three evaluations be read together. The recipe generalizes: if the non-Clifford content of a code sits only in its stabilizers and the state stays in a flat (or otherwise affine) subspace, a code-specific affine-quadratic frame is the fast route; if ideal logical non-Clifford unitaries act on such a code, an intertwiner $C$ with conditions of the type C1--C3 reduces the protocol to a Clifford one; for gate-level noise, Algorithm~2 of Ref.~\cite{takada2026exact} applies once its applicability conditions have been checked.

\paragraph{Limitations.}
All sizes are finite, and the asymptotic thresholds and the fault distance of the full protocol are not determined. The decoders are not optimized. Several elements are idealized: the initial state and termination of the memory experiment, $U$ and the corrections in study 2, feed-forward in study 3, and the boundary qubits fixed by a single $Z$. The CCZ gate is native. The logical states are specific: one holonomy direction in study 1, and input $\ket{+}$ in studies 2 and 3. Phenomenological and circuit-level results use different noise and are not directly comparable.

%=====================================================================
\section{Conclusion}
\label{sec:conclusion}

We evaluated code switching of the $\mathbb{D}_4$ Clifford-hierarchy code at three levels using exact simulations. The $\mathbb{D}_4$ memory has phenomenological finite-size crossings of $2.25$, $2.09$ and $2.00\%$, and its dressed measurements are random both near fluxes and, through CZ-type residuals, without them. The complete switch crosses at $2.23$ and $1.94\%$, well below its idealized version ($2.97$, $2.41\%$), because errors that remain at the switching unitary are amplified. Under circuit-level noise, code switching required at least $1.59$ times the volume of 15-to-1 distillation for an output error of $10^{-4}$ in every condition studied, and $2.30$ times in the single data-certified case. The model supports no size that reaches $10^{-6}$, because of errors in the switching window. These numbers were obtained by choosing, for each question, an exact representation within the Clifford-stabilizer framework: a new twisted-frame representation for the memory, an operator-level Pauli proxy with checkable sufficient conditions for the switch with ideal unitaries, and a circuit-level Pauli proxy for the gate-level resource estimate, cross-validated with zero discrepancies. \blue{[Revise after the $10^5$-shot production of study 1.]} Open directions are the circuit-level version of study 2, decoders and circuits for the switching window, idle noise in the memory experiment, finite-size scaling, and verification for arbitrary input states.

\section*{Acknowledgments}
\blue{[Acknowledgments and funding.]}

\begin{thebibliography}{99}
\bibitem{kobayashi2025clifford}
R. Kobayashi, G. Zhu, and P.-S. Hsin.
\newblock Clifford Hierarchy Stabilizer Codes: Transversal Non-Clifford Gates and Magic States.
\newblock \textit{Phys. Rev. Lett.} \textbf{136}, 250802 (2026); \textit{arXiv:2511.02900 [quant-ph]}.

\bibitem{takada2026exact}
Y. Takada, S. D. Bartlett, and D. J. Williamson.
\newblock Exact efficient simulation of noisy logical magic states using Clifford stabilizers.
\newblock \textit{arXiv:2609.16929 [quant-ph]}, 2026.

\bibitem{haga2026aqis}
I. Haga, S. Nishio, and T. Satoh.
\newblock Threshold Estimation of the $\mathbb{D}_4$ Clifford-Hierarchy Stabilizer Code.
\newblock \blue{AQIS 2026 (submitted) [update the bibliographic status].}

\bibitem{davydova2025knots}
M. Davydova, A. Bauer, J. C. Magdalena de la Fuente, M. Webster, D. J. Williamson, and B. J. Brown.
\newblock Universal fault tolerant quantum computation in 2D without getting tied in knots.
\newblock \textit{Phys. Rev. X} (2026); \textit{arXiv:2503.15751 [quant-ph]}, 2025.

\bibitem{magdalena2026}
J. C. Magdalena de la Fuente, N. Feldman, J. Eisert, and A. Bauer.
\newblock High-threshold decoding of non-Pauli codes for 2D universality.
\newblock \textit{arXiv:2604.02033 [quant-ph]}, 2026.

\bibitem{litinski2019game}
D. Litinski.
\newblock A Game of Surface Codes: Large-Scale Quantum Computing with Lattice Surgery.
\newblock \textit{Quantum} \textbf{3}, 128, 2019.

\bibitem{litinski2019magic}
D. Litinski.
\newblock Magic State Distillation: Not as Costly as You Think.
\newblock \textit{Quantum} \textbf{3}, 205, 2019.

\bibitem{bravyikitaev2005}
S. Bravyi and A. Kitaev.
\newblock Universal quantum computation with ideal Clifford gates and noisy ancillas.
\newblock \textit{Phys. Rev. A} \textbf{71}, 022316, 2005.

\bibitem{bravyikonig2013}
S. Bravyi and R. K\"onig.
\newblock Classification of topologically protected gates for local stabilizer codes.
\newblock \textit{Phys. Rev. Lett.} \textbf{110}, 170503, 2013.

\bibitem{eastin-knill2009}
B. Eastin and E. Knill.
\newblock Restrictions on Transversal Encoded Quantum Gate Sets.
\newblock \textit{Phys. Rev. Lett.} \textbf{102}, 110502, 2009.

\bibitem{iqbal2024nonabelian}
M. Iqbal \textit{et al.}
\newblock Non-Abelian topological order and anyons on a trapped-ion processor.
\newblock \textit{Nature} \textbf{626}, 505--511, 2024.

\bibitem{bombin2018jit}
H. Bombin.
\newblock 2D quantum computation with 3D topological codes.
\newblock \textit{arXiv:1810.09571 [quant-ph]}, 2018.

\bibitem{brown2020jit}
B. J. Brown.
\newblock A fault-tolerant non-Clifford gate for the surface code in two dimensions.
\newblock \textit{Sci. Adv.} \textbf{6}, eaay4929, 2020.

\bibitem{gidney2021stim}
C. Gidney.
\newblock Stim: a fast stabilizer circuit simulator.
\newblock \textit{Quantum} \textbf{5}, 497, 2021.

\bibitem{gottesmanknill1998}
D. Gottesman.
\newblock The Heisenberg representation of quantum computers.
\newblock In \textit{Proceedings of the XXII International Colloquium on Group Theoretical Methods in Physics}, pp. 32--43, 1999.

\bibitem{haenel2026tsim}
R. Haenel, X. Luo, and C. Zhao.
\newblock Tsim: Fast Universal Simulator for Quantum Error Correction.
\newblock \textit{arXiv:2604.01059 [quant-ph]}, 2026.

\bibitem{higgott2025pymatching}
O. Higgott and C. Gidney.
\newblock Sparse Blossom: correcting a million errors per core second with minimum-weight matching.
\newblock \textit{Quantum} \textbf{9}, 1600, 2025.

\bibitem{jax2018github}
J. Bradbury \textit{et al.}
\newblock JAX: composable transformations of Python+NumPy programs.
\newblock \url{http://github.com/jax-ml/jax}, 2018.

\bibitem{agresti1998}
A. Agresti and B. A. Coull.
\newblock Approximate is Better than ``Exact'' for Interval Estimation of Binomial Proportions.
\newblock \textit{The American Statistician} \textbf{52}(2), 119--126, 1998.

\bibitem{howard2021cs}
S. R. Howard, A. Ramdas, J. McAuliffe, and J. Sekhon.
\newblock Time-uniform, nonparametric, nonasymptotic confidence sequences.
\newblock \textit{Ann. Statist.} \textbf{49}(2), 1055--1080, 2021.

\bibitem{ville1939}
J. Ville.
\newblock \textit{\'Etude critique de la notion de collectif}.
\newblock Gauthier-Villars, Paris, 1939.

\end{thebibliography}

\end{document}
